\ifdefined\gkver\else\def\gkver{student}\fi
\documentclass[\gkver]{gaokaozhenti}
\begin{document}

\chapter{2014年重庆理综}
%% paperType: 理综物理部分

\begin{enumerate}
\item
%% number: 1
%% paperName: 2014年重庆理综第 1 题 6 分
%% typeId: 1
%% type: 单选题
%% chapter: 原子和原子核
%% point: 半衰期
%% method: 无
%% score: 6
%% degree: 650
%% duplicateId: 0
%% body:
碘 131 的半衰期约为 8 天，若某药物含有质量为 $m$ 的碘 131，经过 32 天后，该药物中碘 131 的含量大约还有\xzanswer{C}

\fourchoices[answer=C]
{$m/4$}
{$m/8$}
{$m/16$}
{$m/32$}

%% answer: C
%% memo:
\memoanswer{
根据原子核的衰变公式：$m_{\text{剩}}=m\left(\frac{1}{2}\right)^n$，其中 $n=\frac{t}{T}=\frac{32}{8}=4$ 为半衰期的次数，故 $m_{\text{剩}}=\frac{m}{16}$。选项 C 正确。
}

\item
%% number: 2
%% paperName: 2014年重庆理综第 2 题 6 分
%% typeId: 1
%% type: 单选题
%% chapter: 功和能
%% point: 交通工具启动
%% method: 无
%% score: 6
%% degree: 650
%% duplicateId: 0
%% body:
某车以相同功率在两种不同的水平路面上行驶，受到的阻力分别为车重的 $k_{1}$ 和 $k_{2}$ 倍，最大速率分别为 $v_{1}$ 和 $v_{2}$，则\xzanswer{B}

\fourchoices[answer=B]
{$v_{2}=k_{1}v_{1}$}
{$v_{2}=\frac{k_{1}}{k_{2}}v_{1}$}
{$v_{2}=\frac{k_{2}}{k_{1}}v_{1}$}
{$v_{2}=k_{2}v_{1}$}

%% answer: B
%% memo:
\memoanswer{
以相同功率在两种不同的水平路面上行驶达最大速度时，有 $F=F_f=kmg$。由 $P=Fv$ 可知最大速度 $v_m=\frac{P}{F}=\frac{P}{F_f}$，则 $\frac{v_1}{v_2}=\frac{F_{f2}}{F_{f1}}=\frac{k_2}{k_1}$，故 $v_2=\frac{k_1}{k_2}v_1$，选项 B 正确。
}

\item
%% number: 3
%% paperName: 2014年重庆理综第 3 题 6 分
%% typeId: 1
%% type: 单选题
%% chapter: 电场
%% point: 等势面
%% method: 无
%% score: 6
%% degree: 450
%% duplicateId: 0
%% body:
如图所示为某示波管内的聚焦电场，实线和虚线分别表示电场线和等势线。两电子分别从 a、b 两点运动到 c 点，设电场力对两电子做的功分别为 $W_{a}$ 和 $W_{b}$，a、b 两点的电场强度大小分别为 $E_{a}$ 和 $E_{b}$，则\xzanswer{A}

\onepicture[w=5.2cm]{figs/03.png}

\fourchoices[answer=A]
{$W_{a}=W_{b}$，$E_{a}>E_{b}$}
{$W_{a}\neq W_{b}$，$E_{a}>E_{b}$}
{$W_{a}=W_{b}$，$E_{a}<E_{b}$}
{$W_{a}\neq W_{b}$，$E_{a}<E_{b}$}

%% answer: A
%% memo:
\memoanswer{
电子在电场中运动，电场力做功 $W=qU$，a、b 位于同一条等势线上，故 $U_{ac}=U_{bc}$，所以 $W_a=W_b$；而电场线的疏密反应电场强度的大小，a 点比 b 点的电场线密些，故 $E_a>E_b$。选项 A 正确。
}

\item
%% number: 4
%% paperName: 2014年重庆理综第 4 题 6 分
%% typeId: 1
%% type: 单选题
%% chapter: 运动和力
%% point: 动量守恒定律
%% method: 无
%% score: 6
%% degree: 450
%% duplicateId: 0
%% body:
一弹丸在飞行到距离地面 $5\Um$ 高时仅有水平速度 $v=2\Ums$，爆炸成为甲、乙两块水平飞出，甲、乙的质量比为 $3:1$。不计质量损失，取重力加速度 $g=10\Umsq$。则下列图中两块弹片飞行的轨迹可能正确的是\xzanswer{B}

\fourchoices[answer=B, ispicture=true, h=2.2cm]
{figs/04a.png}
{figs/04b.png}
{figs/04c.png}
{figs/04d.png}

%% answer: B
%% memo:
\memoanswer{
弹丸水平飞行爆炸时，在水平方向系统动量守恒，设 $m_{\text{乙}}=m$，则 $m_{\text{甲}}=3m$，故爆炸前水平方向总动量 $p=(3m+m)v=8m$。而爆炸后两弹片做平抛运动，由平抛运动规律：$h=\frac{1}{2}gt^2$，$x_{\text{甲}}=v_{\text{甲}}t$，$x_{\text{乙}}=v_{\text{乙}}t$。

选项 A 中：$v_{\text{甲}}=2.5\Ums$，$v_{\text{乙}}=0.5\Ums$（向左），$p'=3m\times2.5+m\times(-0.5)=7m$，不满足动量守恒，A 错误；

选项 B 中：$v_{\text{甲}}=2.5\Ums$，$v_{\text{乙}}=0.5\Ums$，$p'=3m\times2.5+m\times0.5=8m$，满足动量守恒，B 正确；

选项 C 中：$v_{\text{甲}}=2\Ums$，$v_{\text{乙}}=1\Ums$，$p'=3m\times2+m\times1=7m$，不满足动量守恒，C 错误；

选项 D 中：$v_{\text{甲}}=2\Ums$，$v_{\text{乙}}=1\Ums$（向左），$p'=3m\times2+m\times(-1)=6m$，不满足动量守恒，D 错误。
}

\item
%% number: 5
%% paperName: 2014年重庆理综第 5 题 6 分
%% typeId: 1
%% type: 单选题
%% chapter: 运动和力
%% point: 变加速直线运动
%% method: 无
%% score: 6
%% degree: 450
%% duplicateId: 0
%% body:
以不同初速度将两个物体同时竖直向上抛出并开始计时，一个物体所受空气阻力可忽略，另一个物体所受空气阻力大小与物体速率成正比。下列用虚线和实线描述两物体运动的 $v-t$ 图象可能正确的是\xzanswer{D}

\fourchoices[answer=D, ispicture=true, h=2.4cm]
{figs/05a.png}
{figs/05b.png}
{figs/05c.png}
{figs/05d.png}

%% answer: D
%% memo:
\memoanswer{
竖直上抛运动不受空气阻力，做向上匀减速直线运动至最高点再向下自由落体运动，$v-t$ 图象是倾斜向下的直线，四个选项（虚线）均正确表示；有阻力 $F_f=kv$ 的上抛运动，上升时：$a_1=\frac{mg+kv}{m}$，随着 $v$ 减小，加速度减小，对应的 $v-t$ 图线的斜率减小，A 错误；下落时：$a_2=\frac{mg-kv}{m}$，随着 $v$ 增大，加速度减小，故在最高点时 $v=0$，$a=g$，对应的 $v-t$ 图线与轴的交点，其斜率应该等于 $g$，即过交点的切线应该与竖直上抛运动的直线（虚线）平行，选项 D 正确。
}

\item
%% number: 6
%% paperName: 2014年重庆理综第 6 题 13 分
%% typeId: 4
%% type: 实验题
%% chapter: 力物体的平衡
%% point: 三力平衡
%% method: 无
%% score: 13
%% degree: 450
%% duplicateId: 0
%% body:
\begin{enumerate}
	\item
	某照明电路出现故障，其电路如图 1 所示，该电路用标称值 $12\UV$ 的蓄电池为电源，导线及其接触完好。维修人员使用已调好的多用表直流 $50\UV$ 挡检测故障。他将黑表笔接在 c 点，用红表笔分别探测电路的 a、b 点。

	\twopicture[hA=3.2cm, hB=3.2cm, align=right, numA=1, numB=2]{figs/06a.png}{figs/06b.png}

	\begin{steps}
		\item
		断开开关，红表笔接 $a$ 点时多用表指示如图 2 所示，读数为\tkanswer[1.2cm]{11.5}$\UV$，说明\tkanswer[1.2cm]{蓄电池}正常（选填：蓄电池、保险丝、开关、小灯）。
		\item
		红表笔接 b 点，断开开关时，表针不偏转，闭合开关后，多用表指示仍然和图 2 相同，可判断发生故障的器件是\tkanswer[1.2cm]{小灯}。（选填：蓄电池、保险丝、开关、小灯）
	\end{steps}

	\item
	为了研究人们用绳索跨越山谷过程中绳索拉力的变化规律，同学们设计了如图 1 所示的实验装置。他们将不可伸长轻绳的两端通过测力计（不计质量及长度）固定在相距为 $D$ 的两立柱上，固定点分别为 P 和 Q，P 低于 Q，绳长为 $L$（$L>PQ$）。他们首先在绳上距离 P 点 $10\Ucm$ 处（标记为 C）系上质量为 $m$ 的重物（不滑动），由测力计读出 PC、QC 的拉力大小 $T_{P}$、$T_{Q}$。随后，改变重物悬挂点 C 的位置，每次将 P 到 C 的距离增加 $10\Ucm$，并读出测力计的示数，最后得到 $T_{P}$、$T_{Q}$ 与绳长 PC 的关系曲线如图 2 所示。由实验可知：

	\twopicture[wA=4.5cm, wB=9cm, align=right, numA=1, numB=2]{\includesvg{figs/06c}}{\includesvg{figs/06d}}

	\begin{steps}
		\item
		曲线Ⅱ中拉力最大时，C 与 P 点的距离为\tkanswer[1cm]{60}$\Ucm$，该曲线为\tkanswer[1cm]{$T_{P}$}（选填：$T_{P}$ 或 $T_{Q}$）的曲线。
		\item
		在重物从 P 移到 Q 的整个过程中，受到最大拉力的是\tkanswer[1cm]{Q}（选填：P 或 Q）点所在的立柱。
		\item
		曲线Ⅰ、Ⅱ相交处，可读出绳的拉力为 $T_{0}=$\tkanswer[1cm]{4.30}$\UN$，它与 $L$、$D$、$m$ 和重力加速度 $g$ 的关系为 $T_{0}=$\tkanswer[3cm]{$\frac{mgL}{2\sqrt{L^2-D^2}}$}。
	\end{steps}
\end{enumerate}

%% answer: 
\jdanswer{
\begin{enumerate}
	\item
	① $11.5\UV$（$11.2\sim11.8$ 均可）；蓄电池。② 小灯。

	\item
	① $60\Ucm$（$56\sim64$ 均可），$T_{P}$。② Q。③ $4.30\UN$（$4.25\sim4.35$ 均可），$\frac{mgL}{2\sqrt{L^2-D^2}}$。

\end{enumerate}
}
%% memo:
\memoanswer{
（1）①多用电表直流 $50\UV$ 的量程读中间均匀部分，共 50 格每格 $1\UV$ 应该估读到 $0.1\UV$，指针在 $11\sim12$ 之间，读数为 $11.5\UV$。开关断开，相当于电压表并联在蓄电池两端，读出了蓄电池的电压，故蓄电池是正常的。

②两表笔接 b、c 之间，并闭合开关，测试电压相同，说明 a、b 之间是通路，b、c 之间是断路，故障器件是小灯。

（2）①由曲线Ⅱ的最高点拉力最大，对应的横坐标 $PC=60\Ucm$，设 PC 和 QC 与水平的夹角为 $\alpha$ 和 $\beta$，由 C 点的平衡可得：$T_{P}\cos\alpha=T_{Q}\cos\beta$，开始 C 点靠近 P 点，$\alpha>\beta$，则 $\frac{T_{P}}{T_{Q}}=\frac{\cos\beta}{\cos\alpha}>1$，即 $T_{P}>T_{Q}$，结合两曲线左侧部分，Ⅱ曲线靠上则为 $T_{P}$ 的曲线。

②比较两图象的顶点大小可知，Ⅰ曲线的最高点更大，代表 Q 有最大拉力。

③两曲线的交点表示左右的绳拉力大小相等，读出纵坐标为 $T_P=T_Q=4.30\UN$，设 CQ 绳与立柱的夹角为 $\theta$，延长 CQ 线交另立柱上，构成直角三角形，则 $\cos\theta=\frac{\sqrt{L^2-D^2}}{L}$，由力的平衡可知 $2T_Q\cos\theta=mg$，则 $T_Q=\frac{mg}{2\cos\theta}=\frac{mgL}{2\sqrt{L^2-D^2}}$。
}

\item
%% number: 7
%% paperName: 2014年重庆理综第 7 题 15 分
%% typeId: 6
%% type: 计算题
%% chapter: 功和能
%% point: 动能定理
%% method: 无
%% score: 15
%% degree: 450
%% duplicateId: 0
%% body:
如图所示为“嫦娥三号”探测器在月球上着陆最后阶段的示意图。首先在发动机作用下，探测器受到推力在距月面高度为 $h_{1}$ 处悬停（速度为 0，$h_{1}$ 远小于月球半径）；接着推力改变，探测器开始竖直下降，到达距月面高度为 $h_{2}$ 处的速度为 $v$；此后发动机关闭，探测器仅受重力下落到月面。已知探测器总质量为 $m$（不包括燃料），地球和月球的半径比为 $k_{1}$，质量比为 $k_{2}$，地球表面附近的重力加速度为 $g$。求：

\begin{enumerate}
	\item
	月球表面附近的重力加速度大小及探测器刚接触月面时的速度大小；
	\item
	从开始竖直下降到刚接触月面时，探测器机械能的变化。
\end{enumerate}

\onepicture[w=4cm, align=right]{\includesvg{figs/07}}

%% answer: 
\jdanswer{
\begin{enumerate}
	\item
	$g'=\frac{k_1^2}{k_2}g$，$v_t=\sqrt{v^2+\frac{2k_1^2gh_2}{k_2}}$

	\item
	$\Delta E=\frac{1}{2}mv^2-\frac{k_1^2}{k_2}mg(h_1-h_2)$

\end{enumerate}
}
%% memo:
\memoanswer{
（1）设地球质量和半径分别为 M 和 R，月球的质量、半径和表面附近的重力加速度分别为 $M'$、$R'$ 和 $g'$，探测器刚接触月面时的速度大小为 $v_1$。

由 $mg'=G\frac{M'm}{R'^2}$ 和 $mg=G\frac{Mm}{R^2}$ 得 $g'=\frac{k_1^2}{k_2}g$

由 $v_t^2-v^2=2g'h_2$ 得 $v_t=\sqrt{v^2+\frac{2k_1^2gh_2}{k_2}}$

（2）设机械能变化量为 $\Delta E$，动能变化量为 $\Delta E_k$，重力势能变化量为 $\Delta E_P$。

由 $\Delta E=\Delta E_k+\Delta E_P$

有 $\Delta E=\frac{1}{2}m(v^2+\frac{2k_1^2gh_2}{k_2})+m\frac{k_1^2}{k_2}gh_1$

得：$\Delta E=\frac{1}{2}mv^2-\frac{k_1^2}{k_2}mg(h_1-h_2)$
}

\item
%% number: 8
%% paperName: 2014年重庆理综第 8 题 16 分
%% typeId: 6
%% type: 计算题
%% chapter: 磁场
%% point: 安培力
%% method: 无
%% score: 16
%% degree: 450
%% duplicateId: 0
%% body:
某电子天平原理如图所示，$E$ 形磁铁的两侧为 N 极，中心为 S 极，两极间的磁感应强度大小均为 $B$，磁极宽度均为 $L$，忽略边缘效应。一正方形线圈套于中心磁极，其骨架与秤盘连为一体，线圈两端 C、D 与外电路连接。当质量为 $m$ 的重物放在秤盘上时，弹簧被压缩，秤盘和线圈一起向下运动（骨架与磁极不接触）随后外电路对线圈供电，秤盘和线圈恢复到未放重物时的位置并静止，由此时对应的供电电流 $I$ 可确定重物的质量。已知线圈匝数为 $n$，线圈电阻为 $R$，重力加速度为 $g$。问：

\begin{enumerate}
	\item
	线圈向下运动过程中，线圈中感应电流是从 C 端还是从 D 端流出？
	\item
	供电电流 $I$ 是从 C 端还是从 D 端流入？求重物质量与电流的关系。
	\item
	若线圈消耗的最大功率为 $P$，该电子天平能称量的最大质量是多少？
\end{enumerate}

\onepicture[w=5.5cm, align=right]{figs/08.png}

%% answer: 
\jdanswer{
\begin{enumerate}
	\item
	感应电流从 C 端流出

	\item
	外加电流从 D 端流入，$m=\frac{2nBL}{g}I$

	\item
	$m_0=\frac{2nBL}{g}\sqrt{\frac{P}{R}}$

\end{enumerate}
}
%% memo:
\memoanswer{
（1）感应电流从 C 端流出。

（2）设线圈受到的安培力为 $F_A$，外加电流从 D 点流入。

由 $F_A=mg$ 和 $F_A=2nBIL$ 得：$m=\frac{2nBL}{g}I$

（3）设称量最大质量为 $m_0$

由 $m=\frac{2nBL}{g}I$ 和 $P=I^2R$ 得：$m_0=\frac{2nBL}{g}\sqrt{\frac{P}{R}}$
}

\item
%% number: 9
%% paperName: 2014年重庆理综第 9 题 18 分
%% typeId: 6
%% type: 计算题
%% chapter: 磁场
%% point: 带电粒子在复合场中的运动
%% method: 无
%% score: 18
%% degree: 450
%% duplicateId: 0
%% body:
如图所示，在无限长的竖直边界 NS 和 MT 间充满匀强电场，同时该区域上、下部分分别充满方向垂直于 NSTM 平面向外和向内的匀强磁场，磁感应强度大小分别为 $B$ 和 $2B$，KL 为上下磁场的水平分界线，在 NS 和 MT 边界上，距 KL 高 $h$ 处分别有 P、Q 两点，NS 和 MT 间距为 $1.8h$。质量为 $m$、带电量为 $+q$ 的粒子从 P 点垂直于 NS 边界射入该区域，在两边界之间做圆周运动，重力加速度为 $g$。

\begin{enumerate}
	\item
	求该电场强度的大小和方向；
	\item
	要使粒子不从 NS 边界飞出，求粒子入射速度的最小值；
	\item
	若粒子能经过 Q 点从 MT 边界飞出，求粒子入射速度的所有可能值。
\end{enumerate}

\onepicture[w=5cm, align=right]{figs/09.png}

%% answer: 
\jdanswer{
\begin{enumerate}
	\item
	$E=\frac{mg}{q}$，方向竖直向上；

	\item
	$(9-6\sqrt{2})\frac{Bqh}{m}$；

	\item
	可能的速度有三个：$\frac{0.68Bqh}{m}$，$\frac{0.545Bqh}{m}$，$\frac{0.52Bqh}{m}$。

\end{enumerate}
}
%% memo:
\memoanswer{
（1）设电场强度大小为 $E$，由题意有 $mg=qE$，得 $E=\frac{mg}{q}$，方向竖直向上。

（2）如答题 9 图 1 所示，设粒子不从 NS 边飞出的入射速度最小值为 $v_{\min}$，对应的粒子在上、下区域的运动半径分别为 $r_1$ 和 $r_2$，圆心的连线与 NS 的夹角为 $\varphi$。

由 $r=\frac{mv}{qB}$

有 $r_1=\frac{mv_{\min}}{qB}$，$r_2=\frac{1}{2}r_1$

由 $(r_1+r_2)\sin\varphi=r_2$

$r_1+r_1\cos\varphi=h$

\onepicture[w=4.5cm]{figs/09_an1.jpg}

$v_{\min}=(9-6\sqrt{2})\frac{qBh}{m}$

（3）如答题 9 图 2 所示，设粒子入射速度为 $v$，粒子在上、下方区域的运动半径分别为 $r_1$ 和 $r_2$，粒子第一次通过 KL 时距离 K 点为 $x$，由题意有 $3nx=1.8h(n=1,2,3,\cdots)$

$\frac{3}{2}x\geq\frac{(9-6\sqrt{2})h}{2}$

得 $x=\sqrt{r_1^2-(h-r_1)^2}$

得 $r_1=(1+\frac{0.36}{n^2})\frac{h}{2},\quad n<3.5$

即 $n=1$ 时 $v=\frac{0.68Bqh}{m}$

$n=2$ 时 $v=\frac{0.545Bqh}{m}$

$n=3$ 时 $v=\frac{0.52Bqh}{m}$

\onepicture[w=6.5cm]{figs/09_an2.jpg}
}

\item
%% number: 10
%% paperName: 2014年重庆理综第 10 题 6 分
%% typeId: 6
%% type: 计算题
%% chapter: 气体性质
%% point: 玻意耳定律
%% method: 无
%% score: 6
%% degree: 650
%% duplicateId: 0
%% body:
\begin{enumerate}
	\item
	重庆出租车常以天然气作为燃料，加气站储气罐中天然气的温度随气温升高的过程中，若储气罐内气体体积及质量均不变，则罐内气体（可视为理想气体）\xzanswer{B}

	\fourchoices[answer=B]
	{压强增大，内能减小}
	{吸收热量，内能增大}
	{压强减小，分子平均动能增大}
	{对外做功，分子平均动能减小}

	\item
	如图所示为一种减震垫，上面布满了圆柱状薄膜气泡，每个气泡内充满体积为 $V_{0}$、压强为 $p_{0}$ 的气体。当平板状物品平放在气泡上时，气泡被压缩。若气泡内气体可视为理想气体，其温度保持不变。当体积压缩到 $V$ 时气泡与物品接触面的边界为 $S$。求此时每个气泡内气体对接触面处薄膜的压力。

	\onepicture[w=6.5cm, align=right]{figs/10.png}
\end{enumerate}

%% answer: 
\jdanswer{
\begin{enumerate}
	\item
	B

	\item
	$\frac{V_0}{V}p_0S$

\end{enumerate}
}
%% memo:
\memoanswer{
（1）对理想气体由 $pV=nRT$ 可知体积和质量不变，温度升高时，压强增大，选项 C 错误；理想气体的内能只有分子动能，而温度是分子平均动能的标志，故温度升高，分子平均动能增大，内能增大，选项 A、D 错误；体积不变，故 $W=0$，由热力学第一定律 $\Delta U=Q+W$ 可知，吸热内能增大，选项 B 正确。

（2）设压力为 $F$，压缩后气体压强为 $p$

由等温过程：$p_0V_0=pV$

$F=pS$

解得：$F=\frac{V_0}{V}p_0S$
}

\item
%% number: 11
%% paperName: 2014年重庆理综第 11 题 6 分
%% typeId: 6
%% type: 计算题
%% chapter: 机械振动
%% point: 振动图象
%% method: 无
%% score: 6
%% degree: 650
%% duplicateId: 0
%% body:
\begin{enumerate}
	\item
	打磨某剖面如图所示的宝石时，必须将 OP、OQ 边与轴线的夹角 $\theta$ 切割在 $\theta_{1}<\theta<\theta_{2}$ 的范围内，才能使从 MN 边垂直入射的光线，在 OP 边和 OQ 边都发生全反射（仅考虑如图所示的光线第一次射到 OP 边并反射到 OQ 边后射向 MN 边的情况），则下列判断正确的是\xzanswer{D}

	\onepicture[w=5cm]{figs/11a.png}

	\fourchoices[answer=D]
	{若$\theta>\theta_{2}$，光线一定在 OP 边发生全反射}
	{若$\theta>\theta_{2}$，光线会从 OQ 边射出}
	{若$\theta<\theta_{1}$，光线会从 OQ 边射出}
	{若$\theta<\theta_{1}$，光线会在 OP 边发生全反射}

	\item
	一竖直悬挂的弹簧振子，下端装有一记录笔，在竖直面内放置有一记录纸。当振子上下振动时，以速率 $v$ 水平向左拉动记录纸，记录笔在纸上留下如图所示的图象。$y_{1}$、$y_{2}$、$x_{0}$、$2x_{0}$ 为纸上印迹的位置坐标。由此求振动的周期和振幅。

	\onepicture[w=6.5cm, align=right]{figs/11b.png}
\end{enumerate}

%% answer: 
\jdanswer{
\begin{enumerate}
	\item
	D

	\item
	$T=\frac{2x_0}{v}$，$A=\frac{y_1-y_2}{2}$

\end{enumerate}
}
%% memo:
\memoanswer{
（1）由全反射的临界角满足 $\sin C=1/n$，则入射角满足 $i\geqslant C$ 发生全反射；作出光路可知光线垂直穿过 MN 后到达 OP 的入射角为 $90^\circ-\theta$，则 $\theta$ 越小，越容易发生全反射。选项 D 正确。

（2）设周期为 $T$，振幅为 $A$

由题意得 $T=\frac{2x_0}{v}$ 和 $A=\frac{y_1-y_2}{2}$
}

\end{enumerate}

\end{document}
